\section{User Stories}
\label{sec:userstories}

User Stories beschreiben die Anforderungen aus Sicht der Endbenutzer. Sie folgen dem Format
\textit{``Als [Rolle] moechte ich [Funktion], damit [Nutzen].''} und helfen dabei, den Fokus
auf den tatsaechlichen Mehrwert fuer den Anwender zu legen.

\subsection{Uebersicht der Rollen}

\begin{description}
  \item[Benutzer] Privatperson oder Techniker, der seine Internetqualitaet ueberwachen
    und bei Problemen fundierte Beschwerden beim Provider einreichen moechte.
  \item[Administrator] Person, die SignalReport einrichtet und die Sicherheitseinstellungen
    verwaltet (kann identisch mit dem Benutzer sein).
  \item[System] SignalReport selbst -- fuehrt automatisierte Hintergrundprozesse aus.
\end{description}

\subsection{Kernfunktionen (Must-Have)}

\begin{enumerate}
  \item \textbf{US-01: Kontinuierliche Messung}\\
    \textit{Als Benutzer moechte ich, dass meine Internetverbindung automatisch alle paar
    Sekunden gemessen wird, damit ich lueckenlos dokumentieren kann, wann Probleme auftreten.}
    \begin{itemize}
      \item Akzeptanzkriterien: Ping-, DNS- und HTTP-Messungen laufen im konfigurierbaren
        Intervall (5\,s -- 1\,h); Ergebnisse werden persistent in der H2-Datenbank gespeichert.
    \end{itemize}

  \item \textbf{US-02: Live-Dashboard}\\
    \textit{Als Benutzer moechte ich meine Messergebnisse in Echtzeit im Browser sehen, damit
    ich sofort erkenne, wenn meine Verbindung schlecht ist.}
    \begin{itemize}
      \item Akzeptanzkriterien: Web-UI zeigt Tabelle der letzten Messungen, Live-Chart
        (Chart.js) und Statistik-Panel; Daten aktualisieren sich alle 5 Sekunden.
    \end{itemize}

  \item \textbf{US-03: PDF-Bericht fuer Provider-Beschwerde}\\
    \textit{Als Benutzer moechte ich einen professionellen PDF-Bericht generieren koennen, damit
    ich meinem Provider schwarz auf weiss belegen kann, dass die Verbindung regelmaessig
    einbricht.}
    \begin{itemize}
      \item Akzeptanzkriterien: Berichte fuer 24\,h, 7 Tage und 12 Monate; enthalten
        Kundendaten, Charts mit Zielaenderungsmarkierungen, Top-10 schlechteste Messungen
        und Verbindungsausfaelle.
    \end{itemize}

  \item \textbf{US-04: CSV-Export}\\
    \textit{Als Benutzer moechte ich die Rohdaten als CSV exportieren koennen, damit ich sie in
    Excel oder anderen Tools weiterverarbeiten kann.}
    \begin{itemize}
      \item Akzeptanzkriterien: Vollstaendiger und gefilterter Export moeglich; Semikolon
        als Trennzeichen (fuer deutsche Excel-Versionen).
    \end{itemize}

  \item \textbf{US-05: Statistiken}\\
    \textit{Als Benutzer moechte ich Durchschnittslatenz, 95. Perzentil, Jitter und Paketverlust
    auf einen Blick sehen, damit ich die Qualitaet meiner Verbindung beurteilen kann.}
    \begin{itemize}
      \item Akzeptanzkriterien: Berechnung fuer alle drei Messtypen (PING/DNS/HTTP) ueber
        konfigurierbare Zeitraeume.
    \end{itemize}

  \item \textbf{US-06: Setup-Assistent}\\
    \textit{Als Administrator moechte ich beim ersten Start durch einen Assistenten gefuehrt
    werden, damit ich das System ohne Kommandozeile einrichten kann.}
    \begin{itemize}
      \item Akzeptanzkriterien: Web-basierter Wizard setzt Admin-Passwort und optionale
        Authentifizierung; laeuft nur beim allerersten Start.
    \end{itemize}
\end{enumerate}

\subsection{Erweiterte Funktionen (Should-Have)}

\begin{enumerate}[resume]
  \item \textbf{US-07: DNS-Benchmark}\\
    \textit{Als Benutzer moechte ich verschiedene DNS-Server weltweit vergleichen koennen, damit
    ich den schnellsten DNS-Server fuer meinen Standort finde.}
    \begin{itemize}
      \item Akzeptanzkriterien: Parallele Abfrage von mindestens 10 DNS-Servern aus
        verschiedenen Regionen; Ergebnisse sortiert nach Latenz; 30-Sekunden-Cooldown
        zwischen Benchmarks.
    \end{itemize}

  \item \textbf{US-08: IP-Tracking}\\
    \textit{Als Benutzer moechte ich sehen, wann sich meine externe IP-Adresse aendert, damit
    ich nachvollziehen kann, wie oft mein Provider die Verbindung trennt.}
    \begin{itemize}
      \item Akzeptanzkriterien: Aenderungen werden mit Zeitstempel protokolliert;
        Anzeige im UI mit alter und neuer IP; Statistik ueber Anzahl der Aenderungen.
    \end{itemize}

  \item \textbf{US-09: Wartungsfenster}\\
    \textit{Als Benutzer moechte ich ein Zeitfenster definieren koennen, in dem keine Messungen
    stattfinden, damit geplante Router-Neustarts die Statistik nicht verfaelschen.}
    \begin{itemize}
      \item Akzeptanzkriterien: Start- und Endzeit konfigurierbar; funktioniert auch ueber
        Mitternacht (z.\,B.\ 23:50 -- 00:10); Status im UI sichtbar.
    \end{itemize}

  \item \textbf{US-10: Dark Mode}\\
    \textit{Als Benutzer moechte ich zwischen hellem und dunklem Design wechseln koennen, damit
    die Oberflaeche auch bei naechtlicher Nutzung angenehm ist.}
    \begin{itemize}
      \item Akzeptanzkriterien: Toggle-Button im Header; Einstellung persistent in
        config.json; kein Flackern beim Laden (Server-seitige CSS-Klassen-Injection);
        Chart.js-Farben passen sich an.
    \end{itemize}

  \item \textbf{US-11: Authentifizierung}\\
    \textit{Als Administrator moechte ich den Zugriff auf das Web-Interface mit Passwort
    schuetzen koennen, damit Unbefugte nicht meine Messdaten sehen oder Einstellungen
    aendern koennen.}
    \begin{itemize}
      \item Akzeptanzkriterien: HTTP Basic Auth mit Admin- und User-Rolle; Admin hat
        Vollzugriff, User nur Leserechte; kann jederzeit aktiviert/deaktiviert werden.
    \end{itemize}

  \item \textbf{US-12: Push-Benachrichtigungen}\\
    \textit{Als Benutzer moechte ich bei Verbindungsausfaellen oder hoher Latenz automatisch
    benachrichtigt werden, damit ich auch ohne offenes Browser-Fenster informiert bin.}
    \begin{itemize}
      \item Akzeptanzkriterien: Browser Push Notifications; konfigurierbarer
        Latenz-Schwellwert und Anzahl aufeinanderfolgender schlechter Messungen.
    \end{itemize}
\end{enumerate}

\subsection{Systemfunktionen (automatisiert)}

\begin{enumerate}[resume]
  \item \textbf{US-13: Hintergrund-Dienst}\\
    \textit{Als System moechte ich als Hintergrund-Dienst laufen (auch ohne angemeldeten
    Benutzer), damit die Messung rund um die Uhr erfolgt.}
    \begin{itemize}
      \item Akzeptanzkriterien: Installations-Skripte fuer Windows (procrun),
        Linux (systemd) und macOS (launchd); automatischer Start nach Neustart.
    \end{itemize}

  \item \textbf{US-14: Plattformunabhaengigkeit}\\
    \textit{Als System moechte ich auf Windows, macOS, Linux und Synology NAS lauffaehig sein,
    damit der Benutzer freie Wahl beim Einsatzort hat.}
    \begin{itemize}
      \item Akzeptanzkriterien: Einzelne ausfuehrbare JAR-Datei; embedded H2-Datenbank
        ohne separaten Server; getestet auf allen vier Plattformen.
    \end{itemize}
\end{enumerate}
